\uuid{oUiB}
\titre{Équation droite/plan}
\auteur{bodin ; hanani ; mzari}
\organisation{exo7}
\licence{CC BY-NC-SA 4.0}
\datecreate{2020-01-20}
\niveau{L1}
\module{Géométrie}
\chapitre{Géométrie de l'espace}
\sousChapitre{Distance}
\difficulte{2}
\theme{Distance, équation droite, plan}
\exo7qcmid{lille-1.geometrie-espace.026}

\contenu{
\texte{
Pour ces questions, l'espace est muni d'un repère orthonormé direct  $(O, \vec{i}, \vec{j}, \vec{k})$.
}

\begin{question}
\qclassification{141.05}
\qtags{motcle=équation droite/plan}
On considère les droites $D_1 : \left\{\begin{array}{ccl}x+y-z&=&1\\x-y&=&-1 \end{array}\right.$,
$D_2 : \left\{\begin{array}{ccl}x-y+z&=&-1\\x-z&=&1 \end{array}\right.$
 et $\Delta$ la perpendiculaire commune à $D_1$ et $D_2$. Quelles sont les assertions vraies ?
\begin{answers}

    \bad {Une représentation cartésienne de $\Delta$ est :
 $$\left\{\begin{array}{ccl}x+5y-4z-5&=&0\\x-4y+5z+5&=&0 \end{array}\right.$$}

    \good {Une représentation cartésienne de $\Delta$ est :
 $$\left\{\begin{array}{ccl}x+7y-4z-7&=&0\\x-4y+7z+7&=&0 \end{array}\right.$$}

    \bad { $\Delta$ est contenu dans le plan d'équation $ x+5y-4z-5=0$.}

   \good { $\Delta$ est contenu dans le plan d'équation $x-4y+7z+7=0$.}

\end{answers}
\begin{explanations}
Soit $D_1$ une droite passant par un point $A_1$ et dirigée par un vecteur $\vec{u_1}$ et
 $D_2$ une droite passant par un point $A_2$ et dirigée par un vecteur $\vec{u_2}$, telles que $D_1$ et $D_2$ ne soient pas parallèles. Soit $P_1$ le plan passant par $A_1$ et dirigé par les vecteurs $\vec{u_1}$ et $\vec{u_1} \wedge \vec{u_2}$ et
 $P_2$ le plan passant par $A_2$ et dirigé par les vecteurs $\vec{u_2}$ et $\vec{u_1} \wedge \vec{u_2}$. Alors, la perpendiculaire commune à $D_1$ et $D_2$ est l'intersection de $P_1$ et $P_2$.
\end{explanations}
\end{question}
}
